% ══════════════════════════════════════════════════════════════════
% Part 1: Chapters 1–3
% Spectral Coarse-Graining for Systemic Risk Analysis
% in Financial Networks
% ══════════════════════════════════════════════════════════════════


% ══════════════════════════════════════════════════════════════════
\chapter{Introduction}
\label{ch:introduction}
% ══════════════════════════════════════════════════════════════════

\section{Context and Motivation}
\label{sec:context}

Modern financial systems are characterised by interdependencies that span
multiple temporal, spatial, and relational scales.  At the microscopic
level, individual institutions manage balance sheets, hedge exposures, and
extend credit to counterparties; at the mesoscopic level, clusters of
institutions crystallise around shared asset classes, funding strategies,
or geographic concentration; at the macroscopic level, system-wide
phenomena---contagion cascades, liquidity spirals, fire-sale
externalities---emerge from the collective behaviour of interconnected
agents~\cite{bardoscia2021physics,mantegna1999introduction}.  The global
financial crisis of 2007--2009 demonstrated with devastating clarity that
risk measures evaluated in isolation---Value-at-Risk, stress tests
calibrated to single institutions---fundamentally fail to capture these
emergent dynamics~\cite{brunnermeier2009market,bisias2012survey}.

The crisis exposed what Haldane and May~\cite{haldane2011systemic} termed
the ``robust-yet-fragile'' character of the banking ecosystem: a system
that absorbs small perturbations efficiently through diversification can
nonetheless amplify large shocks catastrophically through the same
channels.  Battiston et al.~\cite{battiston2016price} formalised this
observation as the ``price of complexity,'' showing that the proliferation
of financial instruments and cross-holdings can erode systemic stability
even as individual institutions appear well-capitalised.  These insights
motivate a network perspective: representing institutions as nodes and
their relationships---statistical dependencies, bilateral exposures,
common asset holdings---as edges enables the application of analytical
frameworks from graph theory, statistical physics, and machine
learning~\cite{newman2003structure,barabasi2016network}.

Network representations of financial systems have grown rapidly since
Mantegna's~\cite{mantegna1999hierarchical} pioneering construction of
minimum spanning trees from equity correlation matrices, which revealed
hierarchical taxonomies corresponding to economic sectors.  Subsequent
empirical studies of interbank payment networks~\cite{soramaki2007topology},
overnight lending markets~\cite{iori2008network}, and multiplex exposure
layers~\cite{bargigli2015multiplex,aldasoro2018multiplex} have confirmed
that financial networks exhibit heavy-tailed degree distributions,
small-world properties, and time-varying community structure.  Yet
extracting actionable risk indicators from these networks remains an open
challenge: the high dimensionality and temporal volatility of financial
graphs demand principled methods for dimensionality reduction that
preserve the dynamical properties relevant to contagion and stress
propagation.

This dissertation addresses that challenge by adapting a spectral
coarse-graining methodology---originally developed for general complex
networks---to the domain of interbank systemic risk analysis.  The
approach exploits the eigenstructure of the graph Laplacian to identify
natural scales for network simplification, and integrates the resulting
spectral diagnostics with machine learning and agent-based simulation to
construct a multi-layer analytical framework.

\section{The Spectral Perspective}
\label{sec:spectral_perspective}

The eigenvalues and eigenvectors of the graph Laplacian encode structural
information at every scale of a network~\cite{chung1997spectral}.  The
smallest non-zero eigenvalue, the \emph{algebraic connectivity}
$\lambda_2$ introduced by Fiedler~\cite{fiedler1973algebraic}, governs
the convergence rate of diffusion processes on the graph and determines the
bottleneck resistance to partitioning: a small $\lambda_2$ signals a
network that can be easily split into weakly connected components, while
a large $\lambda_2$ implies rapid mixing and robust
connectivity~\cite{luxburg2007tutorial}.

For financial networks, the spectral perspective offers a natural
language for contagion dynamics.  A shock introduced at a single
institution propagates through the network as a diffusion process
governed by the Laplacian; the spectrum determines the timescales of
this propagation~\cite{masuda2017random}.  Specifically, the
eigenmode expansion of the diffusion operator (Section~\ref{sec:diffusion_theory})
shows that each eigenvalue $\lambda_\alpha$ defines a characteristic
decay time $\tau_\alpha = 1/\lambda_\alpha$: slow modes associated
with small eigenvalues persist and dominate the long-time dynamics,
while fast modes associated with large eigenvalues decay rapidly.

The \emph{spectral gap}---an anomalously large jump between consecutive
eigenvalues---identifies natural timescale separations in the
network~\cite{arenas2006synchronization}.  Arenas, D\'iaz-Guilera, and
P\'erez-Vicente showed that such gaps correspond to the emergence of
community structure at different scales: modes below the gap capture
inter-community dynamics (slow, global), while modes above the gap
capture intra-community dynamics (fast, local).  This timescale
separation is precisely the property that makes principled
coarse-graining possible: if a clear gap exists, one can discard the
fast modes without significantly affecting the long-time behaviour of the
system.

In the financial context, the spectral gap acquires a risk-relevant
interpretation.  A network with a pronounced gap between $\lambda_k$ and
$\lambda_{k+1}$ admits a faithful reduced representation with $k$
effective degrees of freedom.  Our experiments on 1{,}096 daily snapshots
of a 10-bank European interbank network (Section~\ref{ch:results})
confirm that 2--4 clusters emerge across diverse topologies, with the
Erd\H{o}s-R\'enyi gap test rejecting the null hypothesis at $p < 0.01$
in all configurations tested.

\section{Schmidt et al.'s Spectral Coarse-Graining}
\label{sec:intro_scg}

Schmidt, Caccioli, and Aste~\cite{schmidt2024spectral} recently introduced
a spectral coarse-graining and rescaling procedure that generates
reduced-complexity graph representations while preserving both structural
and dynamical properties.  The method draws its conceptual motivation from
the renormalization group in statistical
physics~\cite{wilson1975renormalization}: just as the renormalization group
systematically integrates out short-wavelength fluctuations to produce an
effective theory at longer wavelengths, spectral coarse-graining integrates
out the fast eigenmodes of the graph Laplacian to produce a reduced
network whose diffusion dynamics approximate those of the original system.

The procedure operates in three stages.  First, statistically significant
gaps in the Laplacian spectrum are identified by comparing observed gaps
against an ensemble of Erd\H{o}s-R\'enyi random graphs with matched size
and average degree; gaps whose empirical $p$-value falls below a
threshold $\alpha$ (typically 0.01) are flagged as candidates for
coarse-graining.  Second, the eigenvectors below the gap are used to
partition nodes into clusters via spectral clustering, and a
coarse-grained Laplacian $\mathbf{L}^{(1)}$ is constructed by aggregating
edge weights within and between clusters.  Third, the coarse-grained edge
weights are rescaled by a factor $1/\lambda_k$ to compensate for the
information loss inherent in dimensionality reduction, ensuring that
diffusion dynamics on the reduced graph approximate those of the original
network.

Villegas et al.~\cite{villegas2023laplacian} developed a related
Laplacian renormalization group (LRG) approach based on the normalised
heat kernel.  Schmidt et al.'s method differs in operating directly on
the combinatorial Laplacian spectrum, which better preserves both
dynamics and large-scale topological structure.  The original application
of the method was to electroencephalogram (EEG) recordings of brain
activity, where spectral coarse-graining identified functionally relevant
mesoscale modules in cortical networks.

Adapting this methodology to financial networks is valuable for several
reasons.  First, the principled gap test provides a statistically
grounded criterion for determining the appropriate level of
simplification, unlike ad hoc community detection methods that require
the number of communities as an input parameter.  Second, the rescaling
step preserves diffusion dynamics, which---in the financial
context---correspond to shock propagation processes.  Third, the spectral
quantities computed during the SCG procedure ($\lambda_2$, spectral gap,
spectral radius) are themselves informative risk indicators that can be
tracked over time to monitor changes in network vulnerability.  Our
experiments demonstrate perfect diffusion reconstruction ($R^2 = 1.0$)
across five network topologies, confirming that the method transfers
effectively to financial correlation networks.


\section{Research Objectives and Contributions}
\label{sec:objectives}

This project addresses the gap between spectral graph theory and
operational systemic risk analysis by developing an integrated framework
that combines spectral coarse-graining, graph neural network prediction,
and agent-based stress testing.  The specific research objectives are:

\begin{enumerate}[label=\textbf{O\arabic*.}]
    \item \textbf{Adapt spectral coarse-graining to financial networks.}
    Implement Schmidt et al.'s full renormalization procedure---including
    the Erd\H{o}s-R\'enyi null model gap test with automatic Laplacian
    type detection (combinatorial vs.\ normalised)---for networks
    constructed from daily bank stock return correlations, and evaluate
    preservation of diffusion dynamics across diverse network topologies.

    \item \textbf{Develop a temporal spectral predictor.}  Design and
    train a Graph Attention Network with LSTM temporal
    module~\cite{velickovic2018graph} that predicts the evolution of
    spectral risk indicators ($\lambda_2$, spectral gap, spectral radius)
    from sequences of weekly graph snapshots, enabling forward-looking
    risk assessment.

    \item \textbf{Implement multi-channel agent-based contagion.}
    Construct a three-channel agent-based model capturing credit
    losses~\cite{nier2007network}, funding liquidity
    stress~\cite{gai2010contagion}, and fire-sale asset
    correlation~\cite{caccioli2015overlapping}, and quantify the
    amplification effect relative to a single-channel baseline.

    \item \textbf{Backtest the SCG risk indicator.}  Propose and
    empirically evaluate a spectral risk indicator
    $\text{SCG-Risk} = 1 - \lambda_2/\rho$ as a potential early-warning
    signal for systemic stress, comparing its predictive power against
    realised market volatility and established metrics.

    \item \textbf{Deploy an integrated analytical framework.}  Build a
    modular Python package with an interactive Dash dashboard for
    real-time spectral analysis, GNN training with live progress
    monitoring, and scenario-based stress testing.
\end{enumerate}

The principal contributions of this work are:

\begin{itemize}
    \item The first adaptation of Schmidt et al.'s spectral
    coarse-graining methodology to financial networks, achieving $R^2 =
    1.0$ diffusion reconstruction with statistically significant gap
    detection ($p < 0.01$) across all topologies tested, with 2--4
    clusters identified depending on network construction.

    \item A GAT-LSTM architecture using PyTorch
    Geometric~\cite{fey2019fast} that achieves test $R^2 = 0.382 \pm
    0.072$ for spectral property prediction on 220 weekly snapshots of
    real market data, with per-target analysis revealing algebraic
    connectivity as the most predictable spectral quantity ($R^2 = 0.493$)
    and spectral gap as the least ($R^2 = 0.223$).

    \item Extension of agent-based contagion from a single credit channel
    to a three-channel model (credit, funding liquidity, fire-sale),
    demonstrating $10\times$ default amplification under severe liquidity
    stress (100 Monte Carlo runs per scenario).

    \item Honest empirical evaluation showing that the proposed SCG risk
    indicator exhibits negative rank correlation with future market
    volatility ($r = -0.43$ to $-0.49$), while algebraic connectivity
    alone shows positive predictive power ($r = +0.40$ to $+0.46$),
    revealing a critical distinction between network fragmentation and
    systemic stress.

    \item Cross-correlation analysis exposing the algebraic
    structure of spectral indicators: $\lambda_2$ vs.\ spectral radius
    $\rho = -0.926$, $\lambda_2$ vs.\ SCG risk $= -0.990$, and
    $\lambda_2$ vs.\ realised volatility $= 0.576$.

    \item Interpretable attention analysis revealing that BNP Paribas and
    HSBC are the most attended-to nodes (highest incoming attention),
    while Santander and Deutsche Bank exhibit the highest outgoing
    attention, suggesting that peripheral banks attend most to core
    institutions for spectral prediction.
\end{itemize}


\section{Scope and Limitations}
\label{sec:scope}

Several important limitations bound the scope of this work and should be
stated explicitly.

\paragraph{Network scale.}
The analysis uses a panel of 10 European bank stocks: Deutsche Bank
(DE\_DBK), BNP Paribas (FR\_BNP), Santander (ES\_SAN), UniCredit
(IT\_UCG), ING (NL\_ING), Nordea (SE\_NDA), UBS (CH\_UBS), Barclays
(UK\_BARC), HSBC (UK\_HSBC), and Cr\'edit Agricole (FR\_ACA).  While
this sample covers the major European jurisdictions and represents a
combined balance sheet of several trillion euros, 10 nodes is small for
network-scientific purposes.  Spectral coarse-graining delivers its
greatest value for networks with 50--100+ nodes, where dimensionality
reduction is genuinely necessary; at $N=10$, the SCG procedure
demonstrates technical feasibility but not operational utility.

\paragraph{Correlation proxy.}
The adjacency matrix is constructed from pairwise Pearson correlations of
daily log returns, thresholded to produce binary or weighted networks.
This is a statistical proxy for actual bilateral interbank exposures,
which are not publicly available at the institution level.  Correlation
networks capture co-movement in equity prices, which reflects market
perceptions of interconnectedness rather than the contractual
relationships through which contagion propagates.  The threshold
sensitivity analysis (Section~\ref{ch:results}) reveals that density
drops from 0.956 at threshold 0.2 to 0.473 at threshold 0.6, confirming
that the choice of threshold materially affects network topology and
downstream spectral properties.

\paragraph{Data constraints.}
The dataset comprises 1{,}096 daily snapshots over 2021--2024, yielding
220 weekly observations after stride-5 aggregation.  This provides
sufficient data for training and evaluation, but the period does not
include a major systemic crisis, which limits the ability to validate
early-warning properties of the spectral indicators under genuine stress
conditions.

\paragraph{Indicator inversion.}
The backtesting results reveal that the proposed SCG risk indicator
$1 - \lambda_2/\rho$ is negatively correlated with future market
volatility ($r = -0.43$ to $-0.49$), while volatility persistence
alone achieves $r = 0.67$.  This is an honest negative result: the
indicator captures network fragmentation rather than systemic stress,
and should not be interpreted as an operational early-warning signal
without substantial modification.

\section{Document Structure}
\label{sec:structure}

The remainder of this document is organised as follows. Chapter~\ref{ch:background} reviews the literature on financial networks, contagion, spectral methods, GNNs, ABMs, and risk metrics. Chapter~\ref{ch:theory} develops the mathematical foundations: Laplacian spectral theory, diffusion dynamics, the SCG procedure, GAT attention, and multi-channel contagion. Chapter~\ref{ch:methodology} details the implementation of each component. Chapter~\ref{ch:results} presents experimental results across SCG validation, GAT training, ABM stress testing, backtesting, and sensitivity analysis. Chapter~\ref{ch:discussion} critically analyses the findings and identifies future directions. Chapter~\ref{ch:ethics} addresses ethical considerations. Chapter~\ref{ch:conclusion} summarises the contributions.


% ══════════════════════════════════════════════════════════════════
\chapter{Background and Literature Review}
\label{ch:background}
% ══════════════════════════════════════════════════════════════════

This chapter surveys the six bodies of literature upon which the present
work draws: financial network analysis (Section~\ref{sec:fin_networks}),
systemic risk and contagion mechanisms
(Section~\ref{sec:systemic_risk}), spectral graph theory for financial
networks (Section~\ref{sec:spectral_background}), spectral
coarse-graining and renormalization (Section~\ref{sec:scg_background}),
graph neural networks (Section~\ref{sec:gnn_background}), agent-based
models in financial systems (Section~\ref{sec:abm_background}), and
systemic risk metrics (Section~\ref{sec:risk_metrics_background}).  Each
section critically assesses the strengths and limitations of existing
approaches.  Section~\ref{sec:research_gap} synthesises the identified
gaps and positions this work within the literature.


% ──────────────────────────────────────────────────────────────────
\section{Financial Network Analysis}
\label{sec:fin_networks}
% ──────────────────────────────────────────────────────────────────

\subsection{Network Construction Methodologies}
\label{sec:network_construction}

Three principal approaches to constructing financial networks have
emerged in the literature, each capturing a different facet of
inter-institutional relationships.

\paragraph{Correlation-based networks.}
The econophysics tradition, initiated by Mantegna~\cite{mantegna1999hierarchical},
constructs networks from the pairwise correlation matrix of asset
returns.  Mantegna showed that the minimum spanning tree (MST) of the
correlation distance matrix $d_{ij} = \sqrt{2(1-\rho_{ij})}$ reveals a
hierarchical taxonomy of equities corresponding to industry sectors and
geographic regions~\cite{mantegna1999introduction}.  Tumminello
et al.~\cite{tumminello2005tool} extended this with the Planar Maximally
Filtered Graph (PMFG), which retains substantially more topological
information than the MST while remaining planar, and demonstrated that
the PMFG preserves economically meaningful structures in the correlation
matrix~\cite{tumminello2007correlation}.  Bonanno
et al.~\cite{bonanno2004networks} provided a comprehensive
characterisation of equity correlation networks, establishing that they
exhibit both high clustering and small-world properties.  A limitation
of correlation-based approaches is their sensitivity to the choice of
threshold, estimation window, and frequency: Laloux
et al.~\cite{laloux1999noise} showed via random matrix theory that a
substantial fraction of eigenvalues in empirical correlation matrices are
indistinguishable from noise, implying that naive thresholding can
produce networks dominated by statistical artefacts.

\paragraph{Exposure-based networks.}
Bilateral exposure data---credit claims, derivative positions, securities
cross-holdings---provide a more direct representation of the channels
through which contagion can propagate.  Upper~\cite{nier2007network}
and Nier et al.\ used balance-sheet data to construct interbank lending
networks and simulate cascading defaults, finding that network topology
has a first-order effect on systemic stability.  Elliott, Golub, and
Jackson~\cite{elliott2014financial} modelled cross-holdings as a linear
integration matrix, deriving conditions under which discontinuous
cascades emerge.  However, bilateral exposure data are rarely available
at the institution level: regulators collect such data through exercises
like the European Banking Authority's transparency exercises, but public
releases are aggregated and incomplete.

\paragraph{Payment and transaction networks.}
Soram\"aki et al.~\cite{soramaki2007topology} provided one of the first
comprehensive analyses of interbank payment flows using data from
Fedwire, the large-value real-time gross settlement system operated by
the U.S. Federal Reserve.  They documented small-world properties,
heavy-tailed degree distributions, and a core-periphery structure in
which a small number of highly connected banks intermediate the majority
of payment flows.  Iori et al.~\cite{iori2008network} performed a
parallel analysis of the Italian e-MID overnight lending market,
documenting how network topology evolved through the 2007--2008 crisis:
the network became sparser, more concentrated, and more hierarchical as
banks withdrew from unsecured lending.

\subsection{Network Topology and Stylised Facts}
\label{sec:topology}

Empirical studies have established several stylised facts about financial
network topology.  First, degree distributions are heavy-tailed but
generally not pure power laws; truncated power laws and log-normal
distributions provide better fits~\cite{soramaki2007topology}.  Second,
financial networks exhibit a pronounced core-periphery structure, with a
densely connected core of systemically important institutions and a
sparse periphery of smaller banks~\cite{newman2003structure}.  Third,
financial networks are temporal: their topology evolves substantially
over time, particularly during periods of
stress~\cite{iori2008network,holme2012temporal}.

\subsection{Multiplex Structure}
\label{sec:multiplex}

A critical development in financial network analysis has been the
recognition that interbank relationships form a \emph{multiplex}
structure: different relationship types---unsecured lending, secured
lending, derivatives, securities holdings---constitute distinct network
layers with potentially different topological
properties~\cite{bargigli2015multiplex}.  Bargigli et al.\ analysed the
Italian interbank market and showed that the degree distributions, degree
correlations, and clustering coefficients differ markedly across layers,
implying that single-layer representations lose important structural
information.  Aldasoro and Alves~\cite{aldasoro2018multiplex} confirmed
this finding using European data and demonstrated that systemic
importance rankings of individual institutions depend on which layer (or
combination of layers) is considered.

The multiplex perspective is relevant to the present work because
correlation-based networks---the type we construct from equity return
data---represent a projection of this multiplex structure onto a single
layer.  The correlation between two bank stocks reflects the market's
aggregate assessment of their co-movement, which is influenced by all
underlying exposure channels but does not identify which channels are
active.  This conflation is a fundamental limitation that should be borne
in mind when interpreting spectral properties of correlation networks.


% ──────────────────────────────────────────────────────────────────
\section{Systemic Risk and Contagion Mechanisms}
\label{sec:systemic_risk}
% ──────────────────────────────────────────────────────────────────

\subsection{Direct Contagion: Credit Losses and DebtRank}
\label{sec:direct_contagion}

The most studied contagion mechanism operates through direct bilateral
exposures.  When an institution defaults, its creditors suffer losses
proportional to their exposure and the loss-given-default rate.  Nier
et al.~\cite{nier2007network} showed through simulation that the
relationship between network connectivity and systemic risk is
non-monotonic: moderate connectivity diversifies risk, but high
connectivity facilitates system-wide contagion.

Battiston et al.~\cite{battiston2012debtrank} introduced DebtRank, a
recursive centrality measure that quantifies the fraction of total
economic value in the network that is at risk conditional on the distress
of a particular institution.  Unlike simple cascade models that assume
binary default, DebtRank allows for continuous distress levels that
propagate through the network.  Bardoscia et
al.~\cite{bardoscia2017distress} extended DebtRank to a non-linear
setting in which the propagation of distress depends on the magnitude of
losses, introducing distress propagation functions that capture the
concavity or convexity of the loss-distress relationship.

A key insight from DebtRank analysis is that centrality-based measures of
systemic importance---which depend only on network topology---can diverge
substantially from impact-based measures that account for balance-sheet
heterogeneity~\cite{battiston2009liaisons}.  An institution with
moderate connectivity but large balance sheet may be more systemically
important than a highly connected institution with small exposures.

\subsection{Phase Transitions and the Robust-Yet-Fragile Property}
\label{sec:phase_transitions}

Two influential models characterise the conditions under which contagion
propagates through financial networks.

Elliott, Golub, and Jackson~\cite{elliott2014financial} model an economy
in which organisations hold cross-holdings in one another.  They show
that as the degree of integration increases, the system transitions
discontinuously from a regime in which failures are isolated to one in
which failures cascade through the entire network.  The location of this
phase transition depends on the network topology, the distribution of
initial shocks, and the threshold at which an organisation's value
decline triggers discontinuous failure.

Acemoglu, Ozdaglar, and Tahbaz-Salehi~\cite{acemoglu2015systemic}
formalised the ``robust-yet-fragile'' property more rigorously.  In their
model, a dense (complete) interbank network is optimal for absorbing
small shocks---each institution bears only a small fraction of any
counterparty's losses---but becomes the worst-possible topology for
large shocks, because losses that exceed the absorptive capacity of the
system spread to every institution simultaneously.  The critical shock
size that separates the two regimes decreases with network density,
creating a non-monotonic relationship between connectivity and stability.

These theoretical results have direct implications for spectral
coarse-graining.  The algebraic connectivity $\lambda_2$, which measures
how well-connected the network is, captures the same structural feature
that determines the robust-yet-fragile threshold.  A network with high
$\lambda_2$ diffuses small shocks efficiently but is vulnerable to large
shocks; a network with low $\lambda_2$ may localise large shocks but
fails to absorb even moderate ones.  Our spectral evolution data
(Section~\ref{ch:results}) show $\lambda_2$ with mean $0.923 \pm 0.160$
over the 2021--2024 period, with substantial temporal variation that
tracks changes in market conditions.

\subsection{Indirect Contagion: Funding Liquidity and Fire Sales}
\label{sec:indirect_contagion}

Gai and Kapadia~\cite{gai2010contagion} developed an analytical model of
contagion in random financial networks that emphasises the ``knife-edge''
property: small changes in network structure---a single additional
default, a marginal increase in correlations---can tip the system from
stability to widespread collapse.  Their model captures both direct
credit contagion and the indirect effects of liquidity hoarding.

Brunnermeier and Pedersen~\cite{brunnermeier2009market} provided a
foundational analysis of the feedback loop between market liquidity
(the ease of trading an asset) and funding liquidity (the ease of
obtaining financing).  They showed that a decline in market liquidity
tightens funding constraints, which forces leveraged institutions to
deleverage by selling assets, which further depresses market liquidity.
This ``liquidity spiral'' represents a contagion mechanism that operates
through market prices rather than bilateral exposures.

Caccioli et al.~\cite{caccioli2015overlapping} demonstrated that
overlapping portfolios create a distinct contagion channel.  When
institution $A$ is forced to sell an asset, the price decline imposes
mark-to-market losses on institution $B$ if it holds the same asset,
even if $A$ and $B$ have no direct bilateral relationship.  This
fire-sale externality has been shown to be quantitatively significant:
Caccioli et al.\ documented that fire-sale losses can exceed direct
credit losses in stressed scenarios, particularly in systems with high
leverage~\cite{caccioli2009eroding}.

Glasserman and Young~\cite{glasserman2015contagion} provided analytical
bounds on the probability and magnitude of contagion, showing that
contagion amplification is more severe when exposures are concentrated
rather than diversified.  Their results complement those of Acemoglu et
al.\ by demonstrating that the distribution of exposures, not merely
their mean level, determines systemic vulnerability.

The multi-channel nature of contagion---direct credit, funding liquidity,
and fire-sale externalities---motivates the three-channel agent-based
model developed in this work.  Our ABM experiments
(Section~\ref{ch:results}) show that the funding liquidity channel is
the primary amplification mechanism, producing $10\times$ default
amplification under severe stress compared to a single-channel credit
baseline.

\subsection{The Ecological Analogy}
\label{sec:ecology}

The ecological analogy introduced by May~\cite{may1972will}---that
increased complexity and connectivity in ecosystems can paradoxically
decrease stability---was applied to banking by Haldane and
May~\cite{haldane2011systemic}.  They argued that the banking ecosystem
exhibits a May-Wigner instability threshold, where the product of the
number of institutions, their average connectivity, and the variance of
interaction strengths determines whether the system is stable or unstable.
This insight underscores the importance of spectral methods: the
stability of a linear system is determined by the eigenvalues of its
interaction matrix, and the transition from stability to instability
occurs precisely when the largest eigenvalue crosses a critical
threshold~\cite{may1972will}.

Caccioli et al.~\cite{caccioli2018network} provided a comprehensive
review synthesising network models of systemic risk, while Bardoscia et
al.~\cite{bardoscia2021physics} offered the physics community's
perspective on the statistical mechanics of financial contagion.
Battiston et al.~\cite{battiston2016price} quantified the complexity--stability
trade-off specifically for financial systems.


% ──────────────────────────────────────────────────────────────────
\section{Spectral Graph Theory for Financial Networks}
\label{sec:spectral_background}
% ──────────────────────────────────────────────────────────────────

\subsection{The Laplacian Spectrum}
\label{sec:laplacian_spectrum}

The graph Laplacian and its spectral decomposition form the mathematical
backbone of this work.  Chung's monograph~\cite{chung1997spectral}
provides the foundational treatment, establishing that the eigenvalues
of the normalised Laplacian encode information about expansion,
diameter, chromatic number, and random walk mixing times.  Cvetkovi\'c,
Doob, and Sachs~\cite{cvetkovic1995spectra} developed the broader theory
of graph spectra, covering adjacency, Laplacian, and signless Laplacian
matrices.  Fiedler~\cite{fiedler1973algebraic} introduced the algebraic
connectivity $\lambda_2$ and showed that it provides both upper and
lower bounds on the vertex and edge connectivity of the graph.

For financial networks, the Laplacian spectrum has been used primarily
for two purposes.  First, spectral clustering methods---surveyed by von
Luxburg~\cite{luxburg2007tutorial}---use the bottom eigenvectors of the
normalised Laplacian to partition the network into communities,
following the normalised cut criterion of Ng, Jordan, and
Weiss~\cite{ng2002spectral}.  In the financial context, these
communities correspond to groups of institutions with similar risk
profiles or co-movement patterns.  Second, random matrix theory has
been applied to financial correlation matrices to separate signal from
noise: Laloux et al.~\cite{laloux1999noise} showed that the bulk of
eigenvalues in large correlation matrices follows the Marc\v{e}nko-Pastur
distribution, and only eigenvalues exceeding the upper bound of this
distribution carry genuine information.

\subsection{Spectral Clustering and Community Detection}
\label{sec:spectral_clustering}

Spectral clustering exploits the fact that the first $k$ eigenvectors
of the graph Laplacian define an optimal $k$-dimensional embedding of
the vertices in the sense of minimising the normalised
cut~\cite{ng2002spectral,luxburg2007tutorial}.  In this embedding,
vertices that are densely connected in the original graph are mapped to
nearby points, enabling standard clustering algorithms (typically
$k$-means) to identify communities.  Kyriakopoulos
et al.~\cite{kyriakopoulos2009network} applied spectral methods to
the Greek stock market, using the Laplacian spectrum to identify sector
structure and temporal variation in market organisation.

The relationship between spectral clustering and community detection
methods is mediated by the concept of modularity.  Fortunato and
Hric~\cite{fortunato2016community} provide a comprehensive survey of
community detection methods, including modularity maximisation, label
propagation, and information-theoretic approaches.  Mucha
et al.~\cite{mucha2010community} extended community detection to
temporal and multiplex networks through the concept of multislice
modularity, which is particularly relevant for financial networks that
evolve over time.

\subsection{Temporal Spectral Analysis}
\label{sec:temporal_spectral}

The temporal evolution of spectral properties provides a dynamic view of
network structure that static snapshots cannot capture.  Arenas,
D\'iaz-Guilera, and P\'erez-Vicente~\cite{arenas2006synchronization}
showed that the eigenvalues of the Laplacian determine the
synchronisation properties of coupled oscillators on networks, with
gaps in the spectrum corresponding to timescale separations between
community-level and global dynamics.  Lambiotte, Delvenne, and
Barab\'asi~\cite{lambiotte2014random} connected these spectral
timescales to the Markov time at which random walks on the network
reveal community structure: communities are visible at time $t$ when
$1/\lambda_{k+1} \ll t \ll 1/\lambda_k$ for the gap between eigenvalues
$\lambda_k$ and $\lambda_{k+1}$.

Masuda, Porter, and Lambiotte~\cite{masuda2017random} provide a
comprehensive treatment of random walks on temporal networks, showing
how the interplay between network dynamics and random walk dynamics
gives rise to non-trivial diffusion behaviour.  Holme and
Saram\"aki~\cite{holme2012temporal} survey the broader literature on
temporal networks, including metrics, models, and empirical findings.

Our analysis tracks spectral evolution over 1{,}096 daily snapshots
(Section~\ref{ch:results}), finding that $\lambda_2$ exhibits a mean
of $0.923 \pm 0.160$ with substantial temporal variation, the SCG risk
indicator averages $0.237 \pm 0.181$, and the number of edges varies
between approximately 35 and 46 ($\text{mean} = 40.7 \pm 5.8$),
reflecting changing market conditions over 2021--2024.


% ──────────────────────────────────────────────────────────────────
\section{Spectral Coarse-Graining and Renormalization}
\label{sec:scg_background}
% ──────────────────────────────────────────────────────────────────

\subsection{Renormalization Group Origins}
\label{sec:rg_origins}

The renormalization group (RG), introduced by Wilson~\cite{wilson1975renormalization}
in the context of critical phenomena and quantum field theory, provides a
systematic framework for understanding how physical systems behave across
different scales.  The central idea is to iteratively ``integrate out''
short-wavelength (fast) degrees of freedom, producing an effective theory
for the remaining long-wavelength (slow) degrees of freedom.  At each
step, the coupling constants of the effective theory are
\emph{renormalized} to account for the integrated-out fluctuations.

The application of RG ideas to networks is more recent.  Villegas et
al.~\cite{villegas2023laplacian} developed the Laplacian Renormalization
Group (LRG), which uses the normalised heat kernel at a given diffusion
time $t$ to define a renormalized partition of the network.  At each RG
step, nodes that are ``close'' in diffusion distance at time $t$ are
merged, producing a coarser network.  The procedure can be iterated,
yielding a hierarchy of increasingly coarse representations.  The LRG
is elegant and principled, but operates on the normalised Laplacian and
requires the diffusion time as a parameter, which introduces a
user-specified scale.

\subsection{Schmidt et al.'s Spectral Coarse-Graining}
\label{sec:schmidt_method}

Schmidt, Caccioli, and Aste~\cite{schmidt2024spectral} introduced an
alternative spectral coarse-graining procedure that avoids the
explicit choice of diffusion time.  Instead, the method identifies
natural scales directly from the Laplacian spectrum through a
statistical gap test.  The procedure has three components:

\paragraph{Gap identification.}  For each gap
$\delta_k = |\lambda_{k+1} - \lambda_k|$ in the ordered spectrum,
the observed value is compared against the distribution of gaps from
an ensemble of $10^4$ connected Erd\H{o}s-R\'enyi graphs
$G_{\text{ER}}(n,p)$ with matching number of vertices and expected
average degree.  Gaps whose empirical one-tailed $p$-value falls below
$\alpha = 0.01$ are flagged as statistically significant, identifying
natural scales for coarse-graining.

\paragraph{Spectral clustering and coarse-graining.}  Given a
significant gap at position $k$, the $k$ eigenvectors corresponding
to eigenvalues below the gap are used as features for $k$-means
clustering, partitioning the $N$ nodes into $k$ groups.  The
coarse-grained Laplacian $\mathbf{L}^{(1)} \in \mathbb{R}^{k \times k}$
is then constructed by aggregating edge weights: the off-diagonal
entry $L^{(1)}_{ab}$ equals the sum of original edge weights between
clusters $a$ and $b$, with self-loops accounting for intra-cluster
edges.

\paragraph{Rescaling.}  The coarse-grained adjacency is rescaled by
a factor $1/\lambda_k$ (the eigenvalue at the gap), producing a
rescaled adjacency matrix $\mathbf{A}_R$ whose diffusion dynamics
approximate those of the original network.  This rescaling compensates
for the compression inherent in coarse-graining: without it, the
coarse-grained network would diffuse too slowly because the fast modes
have been removed.

\subsection{Comparison with Alternative Approaches}
\label{sec:scg_comparison}

Schmidt et al.'s method differs from the LRG of Villegas et
al.~\cite{villegas2023laplacian} in three respects.  First, it operates
on the combinatorial Laplacian rather than the normalised Laplacian,
which better preserves the absolute timescales of diffusion processes.
Second, the gap test provides an automatic, statistically grounded
criterion for the coarse-graining resolution, whereas the LRG requires
the diffusion time as an input.  Third, the rescaling step explicitly
compensates for information loss, ensuring that the coarse-grained
dynamics approximate the original dynamics quantitatively.

The method also differs from standard community detection algorithms,
which seek to maximise modularity~\cite{fortunato2016community} or
minimise normalised cuts~\cite{ng2002spectral,luxburg2007tutorial}.
While the spectral clustering step used in SCG is similar to normalised-cut
partitioning, the gap test and rescaling steps are distinct: the gap test
determines \emph{whether} coarse-graining is appropriate (not just how
many communities exist), and the rescaling step ensures dynamical
fidelity rather than merely topological similarity.  Mucha et
al.'s~\cite{mucha2010community} multislice modularity framework provides
a complementary perspective for temporal networks but does not address
the dynamical preservation that is central to SCG.


% ──────────────────────────────────────────────────────────────────
\section{Graph Neural Networks}
\label{sec:gnn_background}
% ──────────────────────────────────────────────────────────────────

\subsection{Historical Development}
\label{sec:gnn_history}

The extension of neural networks to graph-structured data has proceeded
through several conceptual stages.  Scarselli et
al.~\cite{scarselli2009graph} introduced the original Graph Neural
Network (GNN) as a recursive message-passing system in which node
representations are iteratively updated based on their neighbours'
representations until a fixed point is reached.  While theoretically
appealing, this approach suffered from computational limitations and
convergence issues.

The spectral perspective on graph convolution, initiated by Bruna et
al.~\cite{bruna2014spectral}, defined graph convolution in the spectral
domain via the graph Fourier transform---i.e., multiplication in the
eigenspace of the graph Laplacian.  Defferrard et
al.~\cite{defferrard2016convolutional} made this practical by
approximating spectral filters with Chebyshev polynomials, avoiding the
$O(N^3)$ eigendecomposition.  Kipf and
Welling~\cite{kipf2017semi} simplified this further with the Graph
Convolutional Network (GCN), which uses a first-order Chebyshev
approximation with renormalisation.

\subsection{Message Passing and the Spatial Perspective}
\label{sec:mpnn}

Gilmer et al.~\cite{gilmer2017neural} unified the various GNN
architectures under the Message Passing Neural Network (MPNN) framework,
showing that GCNs, GATs, and other variants can all be expressed as
instances of a general message-passing scheme:
\begin{equation}
\mathbf{h}_i^{(\ell+1)} = \phi\!\left(\mathbf{h}_i^{(\ell)},\;
\bigoplus_{j \in \mathcal{N}(i)} \psi\!\left(\mathbf{h}_i^{(\ell)},\,
\mathbf{h}_j^{(\ell)},\, \mathbf{e}_{ij}\right)\right)
\label{eq:mpnn}
\end{equation}
where $\psi$ is the message function, $\bigoplus$ is a permutation-invariant
aggregation (sum, mean, or max), and $\phi$ is the update function.  This
framework clarifies the design space and has guided subsequent
architectural innovation.  Wu et al.~\cite{wu2020comprehensive} provide
a comprehensive survey of the GNN landscape, covering architectures,
training procedures, and applications.

\subsection{Graph Attention Networks}
\label{sec:gat_background}

Veli\v{c}kovi\'c et al.~\cite{velickovic2018graph} introduced Graph
Attention Networks (GATs), which replace the fixed aggregation weights
of GCNs with learned attention coefficients:
\begin{equation}
\alpha_{ij} = \frac{
  \exp\!\left(\text{LeakyReLU}\!\left(\mathbf{a}^T\!
  \left[\mathbf{W}\mathbf{h}_i \,\|\, \mathbf{W}\mathbf{h}_j\right]
  \right)\right)}{
  \sum_{k \in \mathcal{N}(i)}
  \exp\!\left(\text{LeakyReLU}\!\left(\mathbf{a}^T\!
  \left[\mathbf{W}\mathbf{h}_i \,\|\, \mathbf{W}\mathbf{h}_k\right]
  \right)\right)}
\label{eq:gat_attn}
\end{equation}
The attention mechanism offers two advantages over fixed spectral
filters.  First, it is \emph{adaptive}: the weights $\alpha_{ij}$
depend on the current node features, so the aggregation pattern changes
as the input changes.  In financial networks, where the relative
importance of inter-bank relationships varies with market conditions,
this adaptivity is valuable.  Second, attention weights are
\emph{interpretable}: examining which banks attend to which reveals
the network relationships that the model considers most informative for
prediction.

Multi-head attention with $K$ heads stabilises learning and increases
representational capacity~\cite{velickovic2018graph}:
\begin{equation}
\mathbf{h}_i^{(\ell+1)} = \Big\|_{k=1}^{K}\,\sigma\!\left(
\sum_{j \in \mathcal{N}(i)} \alpha_{ij}^{(k)}\,
\mathbf{W}^{(k)}\,\mathbf{h}_j^{(\ell)}\right)
\label{eq:multihead_attn}
\end{equation}
Our implementation uses $K=8$ attention heads with hidden dimension
$h=64$ (Section~\ref{ch:methodology}), yielding 73{,}795 trainable
parameters.

\subsection{Oversmoothing and Depth Limitations}
\label{sec:oversmoothing}

A well-documented limitation of GNNs is \emph{oversmoothing}: as the
number of layers $L$ increases, node representations converge to a
common vector, losing the discriminative power needed for node-level
tasks~\cite{kipf2017semi}.  This is a spectral phenomenon---repeated
application of the graph convolution operator is equivalent to
low-pass filtering, which progressively attenuates the higher-frequency
components that distinguish different nodes.  For small networks like
our 10-bank system, $L=2$ layers suffice to propagate information
across the entire graph (the diameter is at most 2--3 for the dense
topologies we consider), so oversmoothing is not a practical concern.
However, it would become relevant for larger networks where deeper
architectures might be needed.

\subsection{GNNs for Financial Applications}
\label{sec:gnn_finance}

GNNs have been applied to several financial prediction tasks, including
credit scoring, fraud detection, and portfolio
construction~\cite{wu2020comprehensive}.  However, the specific problem
of predicting \emph{spectral properties} of time-evolving financial
networks from graph sequences has not, to our knowledge, been addressed
in the literature.  This is a distinctive feature of the present work:
rather than predicting node-level quantities (e.g., individual bank
returns) or graph-level labels (e.g., crisis / no-crisis), we predict
continuous spectral properties ($\lambda_2$, spectral gap, spectral
radius) that characterise the network as a whole.  The implementation
uses PyTorch Geometric~\cite{fey2019fast}.


% ──────────────────────────────────────────────────────────────────
\section{Agent-Based Models in Financial Systems}
\label{sec:abm_background}
% ──────────────────────────────────────────────────────────────────

\subsection{Motivation and Advantages}
\label{sec:abm_motivation}

Agent-based models (ABMs) represent each institution as an autonomous
agent with state variables (balance sheet, regulatory ratios), decision
rules (lending, borrowing, deleveraging), and interaction mechanisms
(bilateral contracts, market clearing)~\cite{tesfatsion2006agent}.  The
aggregate behaviour of the system emerges from the interactions of
heterogeneous agents, without imposing equilibrium conditions or
representative-agent assumptions.

Farmer and Foley~\cite{farmer2009economy} argued that ABMs offer
advantages over dynamic stochastic general equilibrium (DSGE) models
for studying financial crises: they can incorporate heterogeneous agents,
bounded rationality, network structure, and non-linear feedback loops
that give rise to the fat tails, clustering, and regime shifts observed
in financial data.  Schweitzer et al.~\cite{schweitzer2009economic}
placed this argument in the broader context of complexity economics,
emphasising that emergent phenomena---which are precisely what systemic
risk embodies---require bottom-up modelling approaches.

\subsection{Contagion Models}
\label{sec:abm_contagion}

The seminal financial ABM of Nier et al.~\cite{nier2007network} modelled
banks with stylised balance sheets and simulated cascading defaults on
random interbank networks, establishing that network topology---not just
aggregate capital levels---determines systemic stability.  Gai and
Kapadia~\cite{gai2010contagion} extended this framework analytically,
deriving conditions under which the ``knife-edge'' property arises in
random financial networks.

Multi-channel contagion has been a major theme in recent ABM development.
Poledna et al.~\cite{poledna2015economics} constructed a multi-layer
interbank network model using Austrian supervisory data, demonstrating
that single-layer models substantially underestimate cascading losses
because they miss the amplification from cross-layer interactions.
Montagna and Kok~\cite{montagna2016multi} developed a multi-layered
model for ECB macro-prudential stress testing, incorporating credit,
funding, and market risk channels.  Thurner, Farmer, and
Geanakoplos~\cite{thurner2012leverage} demonstrated the importance of
leverage dynamics: procyclical leverage creates a feedback loop in which
declining asset prices force deleveraging, which further depresses prices.

\subsection{Calibration Challenges}
\label{sec:abm_calibration}

A persistent challenge for financial ABMs is calibration: mapping the
model parameters to empirical data in a way that produces realistic
dynamics.  Unlike econometric models, ABMs typically do not admit
closed-form likelihoods, making standard estimation techniques
inapplicable.  Approximate Bayesian computation, moment matching, and
other simulation-based inference methods have been
proposed~\cite{farmer2009economy}, but remain computationally expensive
and sensitive to the choice of summary statistics.

Our ABM makes no claim to calibration against a specific banking system.
Instead, it uses stylised parameters informed by the literature and
focuses on \emph{relative} comparisons: how much amplification does
three-channel contagion produce relative to single-channel, and how does
the spectral structure of the network modulate this amplification?  This
comparative approach avoids the calibration problem while still yielding
substantive insights into the role of multi-channel dynamics.


% ──────────────────────────────────────────────────────────────────
\section{Systemic Risk Metrics}
\label{sec:risk_metrics_background}
% ──────────────────────────────────────────────────────────────────

\subsection{Market-Based Metrics}
\label{sec:market_metrics}

Several market-based systemic risk metrics have been developed to
complement network-based approaches.

\textbf{SRISK.}  Brownlees and Engle~\cite{brownlees2017srisk} defined
SRISK as the expected capital shortfall of a financial institution
conditional on a prolonged market decline.  The measure combines
market capitalisation, book debt, and the long-run marginal expected
shortfall (LRMES), estimated from an asymmetric dynamic conditional
correlation model.  SRISK has the advantage of being computable from
publicly available data and has been shown to predict the institutions
that required government support during the 2008 crisis.

\textbf{Marginal Expected Shortfall.}  Acharya et
al.~\cite{acharya2017measuring} proposed MES as a measure of the
contribution of an individual institution to aggregate tail risk.  MES
is defined as the expected loss of institution $i$ conditional on the
market being in the $\alpha$-tail of its return distribution.  The
measure captures the co-dependence between individual and system-wide
losses but does not account for network topology or contagion dynamics.

\textbf{Limitations.}  Both SRISK and MES are backward-looking
indicators that rely on historical return distributions and may not
anticipate structural shifts in the network.  Furthermore, they treat
the financial system as a set of return series rather than as a network,
missing the topological information that determines how shocks propagate.

\subsection{Macroprudential Indicators}
\label{sec:macropru}

The macroprudential literature has developed indicators aimed at
identifying the build-up of systemic vulnerabilities before they
materialise.  Drehmann and Juselius~\cite{drehmann2013model} showed
that the credit-to-GDP gap is one of the best single predictors of
banking crises at horizons of 1--3 years.  Borio~\cite{borio2011macroprudential}
articulated the macroprudential policy framework, distinguishing between
the cross-sectional dimension (how risk is distributed across
institutions at a given point in time) and the time-series dimension
(how aggregate risk evolves through the financial cycle).

Admati and Hellwig~\cite{admati2013fallacies} argued that capital
requirements under the Basel framework are insufficient because they
do not account for the systemic externalities of interconnectedness.
Their critique motivates the development of network-based risk
indicators that complement institution-level capital ratios.

\subsection{Survey of the Landscape}
\label{sec:survey}

Bisias et al.~\cite{bisias2012survey} provide a comprehensive survey of
31 quantitative measures of systemic risk, categorising them into
macroeconomic measures, granular measures (including network-based
approaches), and market-based measures.  They observe that no single
measure captures all dimensions of systemic risk, and advocate for a
``dashboard'' approach that monitors multiple indicators simultaneously.
This observation motivates the integrated framework developed in this
work, which combines spectral network diagnostics (SCG), predictive
modelling (GAT-LSTM), and stress simulation (ABM) in a unified
dashboard.


% ──────────────────────────────────────────────────────────────────
\section{Research Gap and Positioning}
\label{sec:research_gap}
% ──────────────────────────────────────────────────────────────────

The preceding review identifies several gaps in the existing literature
that this work aims to address.

\paragraph{Gap 1: Spectral coarse-graining for financial networks.}
While spectral methods have been applied to financial networks for
clustering~\cite{luxburg2007tutorial,kyriakopoulos2009network}, noise
filtering~\cite{laloux1999noise}, and static structural
analysis~\cite{bonanno2004networks}, the application of spectral
\emph{coarse-graining}---as a rigorous dimensionality reduction that
preserves dynamical properties---has not been explored.  Schmidt et
al.'s~\cite{schmidt2024spectral} method has been applied only to EEG
data; its adaptation to financial networks requires addressing the
specific challenges of correlation-based adjacency, threshold
sensitivity, and temporal variability.

\paragraph{Gap 2: GNN prediction of spectral properties.}
Graph neural networks have been applied to financial prediction tasks
such as credit scoring, fraud detection, and return
forecasting~\cite{wu2020comprehensive}, but the specific problem of
predicting spectral properties (algebraic connectivity, spectral gap,
spectral radius) of time-evolving financial networks from graph
sequences has not been addressed.  This is a graph-level regression
task that requires temporal modelling of network evolution, combining
spatial (graph) and temporal (sequence) information.

\paragraph{Gap 3: Integration of spectral diagnostics with ABM stress
testing.}
Agent-based models of financial contagion typically use exogenous network
structures~\cite{nier2007network,poledna2015economics} or stylised
topologies~\cite{gai2010contagion} rather than empirically estimated
networks with spectral characterisation.  Conversely, spectral analyses
of financial networks~\cite{kyriakopoulos2009network} typically stop at
structural description without connecting to contagion simulation.  No
prior work, to our knowledge, integrates SCG-derived spectral
diagnostics with multi-channel ABM stress testing in a unified framework.

\paragraph{Gap 4: Empirical validation of spectral risk indicators.}
The construction of risk indicators from spectral properties (e.g.,
$1 - \lambda_2/\rho$) has theoretical appeal but requires empirical
backtesting against realised outcomes.  The literature on spectral
methods in finance has focused on methodology rather than operational
validation.  Our backtesting analysis (Section~\ref{ch:results})
provides this validation, with the honest finding that the naive SCG
risk indicator fails as an early-warning signal---a result that is
itself informative for future research.

\paragraph{Positioning.}  This work bridges spectral graph theory,
graph neural networks, and agent-based modelling for systemic risk
analysis.  It is the first to adapt Schmidt et al.'s spectral
coarse-graining to financial networks, the first to train a GNN to
predict spectral properties of financial graph sequences, and the
first to integrate spectral diagnostics with multi-channel contagion
simulation.  The framework is implemented as a modular Python package
with an interactive dashboard, enabling reproducible analysis and
extension to larger networks.


% ══════════════════════════════════════════════════════════════════
\chapter{Theoretical Foundations}
\label{ch:theory}
% ══════════════════════════════════════════════════════════════════

This chapter develops the mathematical framework underlying the three
components of our analytical pipeline: spectral coarse-graining
(Sections~\ref{sec:laplacian_theory}--\ref{sec:scg_theory}), graph
attention prediction (Section~\ref{sec:gat_theory}), and multi-channel
contagion simulation (Section~\ref{sec:contagion_theory}).
Section~\ref{sec:risk_formulation} formulates the proposed risk
indicator and discusses its theoretical basis and limitations.


% ──────────────────────────────────────────────────────────────────
\section{Graph Laplacian and Spectral Decomposition}
\label{sec:laplacian_theory}
% ──────────────────────────────────────────────────────────────────

\begin{definition}[Financial Network Graph]
\label{def:financial_graph}
A \emph{financial network graph} is a weighted undirected graph
$\mathcal{G} = (V, E, w)$ where $V = \{v_1, \ldots, v_N\}$ is a set
of $N$ financial institutions, $E \subseteq V \times V$ is the edge
set, and $w : E \to \mathbb{R}_{>0}$ assigns positive weights to
edges.  The adjacency matrix $\mathbf{A} \in \mathbb{R}^{N \times N}$
has entries $A_{ij} = w(v_i, v_j)$ if $(v_i, v_j) \in E$ and
$A_{ij} = 0$ otherwise.
\end{definition}

In this work, $N = 10$ and the weights are derived from pairwise
correlations of daily log returns (Section~\ref{ch:methodology}).  We
construct both binary networks (thresholded at a cutoff $\theta$) and
weighted networks (using the absolute correlation or the correlation
distance $d_{ij} = \sqrt{2(1 - \rho_{ij})}$ as
weight)~\cite{mantegna1999introduction}.

\begin{definition}[Combinatorial and Normalised Laplacian]
\label{def:laplacian}
Given the adjacency matrix $\mathbf{A}$ and the degree matrix
$\mathbf{D} = \operatorname{diag}(d_1, \ldots, d_N)$ with
$d_i = \sum_{j=1}^{N} A_{ij}$, the \emph{combinatorial Laplacian} is
\begin{equation}
\mathbf{L} = \mathbf{D} - \mathbf{A},
\label{eq:comb_laplacian}
\end{equation}
and the \emph{normalised Laplacian} is
\begin{equation}
\mathcal{L} = \mathbf{D}^{-1/2}\,\mathbf{L}\,\mathbf{D}^{-1/2}
= \mathbf{I} - \mathbf{D}^{-1/2}\,\mathbf{A}\,\mathbf{D}^{-1/2}.
\label{eq:norm_laplacian}
\end{equation}
Both matrices are real and symmetric~\cite{chung1997spectral}.
\end{definition}

\begin{theorem}[Spectral Properties of the Laplacian]
\label{thm:spectral}
Let $\mathbf{L}$ be the combinatorial Laplacian of a weighted undirected
graph $\mathcal{G}$ with $N$ vertices.  Then:
\begin{enumerate}[label=(\roman*)]
    \item $\mathbf{L}$ is positive semidefinite: all eigenvalues satisfy
    $\lambda_\alpha \geq 0$.
    \item The smallest eigenvalue is $\lambda_0 = 0$ with eigenvector
    $\mathbf{u}_0 = \frac{1}{\sqrt{N}}\,\mathbf{1}$.
    \item The multiplicity of the zero eigenvalue equals the number of
    connected components of $\mathcal{G}$.
    \item For a connected graph, the eigenvalues satisfy
    $0 = \lambda_0 < \lambda_1 \leq \lambda_2 \leq \cdots \leq
    \lambda_{N-1} \leq 2\,d_{\max}$, where $d_{\max}$ is the maximum
    degree.
\end{enumerate}
\end{theorem}

\begin{proof}[Proof sketch]
For (i), observe that for any vector $\mathbf{x} \in \mathbb{R}^N$,
\begin{equation}
\mathbf{x}^T \mathbf{L}\, \mathbf{x}
= \mathbf{x}^T (\mathbf{D} - \mathbf{A})\, \mathbf{x}
= \sum_{(i,j) \in E} A_{ij}\,(x_i - x_j)^2
\geq 0,
\label{eq:psd_proof}
\end{equation}
which establishes positive semidefiniteness via the quadratic form
characterisation~\cite{chung1997spectral}.  For (ii), $\mathbf{L}\,
\mathbf{1} = (\mathbf{D} - \mathbf{A})\,\mathbf{1} = \mathbf{d} -
\mathbf{d} = \mathbf{0}$ since each row of $\mathbf{L}$ sums to zero.
For (iii), the null space of $\mathbf{L}$ is spanned by indicator
vectors of the connected components~\cite{cvetkovic1995spectra}.  For
(iv), the upper bound follows from the Gershgorin circle
theorem~\cite{golub2013matrix}.
\end{proof}

\begin{definition}[Algebraic Connectivity and Spectral Gap]
\label{def:algebraic_connectivity}
The \emph{algebraic connectivity} of a connected graph $\mathcal{G}$ is
the second-smallest eigenvalue $\lambda_1$ of its Laplacian, also called
the \emph{Fiedler value}~\cite{fiedler1973algebraic}.  The corresponding
eigenvector $\mathbf{u}_1$ is the \emph{Fiedler vector}, whose sign
pattern defines the optimal bisection of the graph.

The \emph{spectral gap} at index $k$ is $\delta_k = \lambda_{k+1} -
\lambda_k$.  A ``large'' spectral gap indicates a natural timescale
separation: the $k$ slowest eigenmodes are well-separated from the
remaining faster modes.
\end{definition}

Note: following convention in the graph Laplacian literature, we index
eigenvalues starting from $\lambda_0 = 0$, so the algebraic
connectivity is $\lambda_1$ in this indexing.  However, when describing
experimental results, we often use $\lambda_2$ to refer to the algebraic
connectivity, following the convention of Schmidt et
al.~\cite{schmidt2024spectral} who index from $\lambda_1$.  Both
conventions appear in the literature; context disambiguates.


% ──────────────────────────────────────────────────────────────────
\section{Diffusion Dynamics on Financial Networks}
\label{sec:diffusion_theory}
% ──────────────────────────────────────────────────────────────────

The graph Laplacian governs continuous-time diffusion on
networks~\cite{masuda2017random}.  Let $\mathbf{p}(t) \in \mathbb{R}^N$
represent the distribution of a quantity (e.g., financial distress)
across $N$ institutions at time $t$.  The diffusion equation is:
\begin{equation}
\frac{d\mathbf{p}}{dt} = -\mathbf{L}\,\mathbf{p}(t).
\label{eq:diffusion_eq}
\end{equation}
This has the formal solution:
\begin{equation}
\mathbf{p}(t) = e^{-t\mathbf{L}}\,\mathbf{p}(0).
\label{eq:diffusion_matrix_exp}
\end{equation}
Expanding in the eigenbasis $\{\mathbf{u}_\alpha\}_{\alpha=0}^{N-1}$ of
$\mathbf{L}$:
\begin{equation}
\mathbf{p}(t) = \sum_{\alpha=0}^{N-1} e^{-t\lambda_\alpha}\,
c_\alpha\,\mathbf{u}_\alpha,
\quad\text{where}\quad
c_\alpha = \mathbf{u}_\alpha^T\,\mathbf{p}(0).
\label{eq:eigenmode_expansion}
\end{equation}

\begin{theorem}[Convergence Rate]
\label{thm:convergence}
For a connected graph with algebraic connectivity $\lambda_1 > 0$, the
solution~\eqref{eq:eigenmode_expansion} converges to the stationary
distribution $\mathbf{p}_\infty = c_0\,\mathbf{u}_0$ at rate
$\lambda_1$:
\begin{equation}
\left\|\mathbf{p}(t) - \mathbf{p}_\infty\right\|_2
\leq \left\|\mathbf{p}(0)\right\|_2\,e^{-\lambda_1 t}.
\label{eq:convergence_bound}
\end{equation}
The characteristic timescale of convergence is
$\tau_{\text{mix}} = 1/\lambda_1$.
\end{theorem}

\begin{proof}[Proof sketch]
Since $\lambda_0 = 0$, the mode $\alpha = 0$ contributes a constant
$c_0\,\mathbf{u}_0$ that persists indefinitely.  All other modes
$\alpha \geq 1$ decay as $e^{-t\lambda_\alpha}$, and the slowest-decaying
non-trivial mode has rate $\lambda_1$.  The bound follows by bounding
$e^{-t\lambda_\alpha} \leq e^{-t\lambda_1}$ for $\alpha \geq
1$~\cite{chung1997spectral}.
\end{proof}

\paragraph{Financial interpretation.}
In the financial context, $\mathbf{p}(0)$ represents an initial shock
localised at one or more institutions, and $\mathbf{p}(t)$ describes how
that shock diffuses through the interbank network.  The algebraic
connectivity $\lambda_1$ determines the \emph{speed} of shock
propagation: a large $\lambda_1$ implies rapid diffusion (the shock
reaches all institutions quickly), while a small $\lambda_1$ implies
slow propagation (the shock may remain localised for an extended
period)~\cite{haldane2011systemic,elliott2014financial}.

The eigenmode expansion~\eqref{eq:eigenmode_expansion} decomposes the
diffusion process into independent modes, each with its own timescale.
The spectral gap identifies the point at which this timescale separation
is most pronounced.  If $\delta_k = \lambda_{k+1} - \lambda_k$ is
large, then the first $k$ modes (slow, global) are well-separated from
the remaining $N - k - 1$ modes (fast, local), and coarse-graining to
$k$ effective degrees of freedom loses little dynamical information.
This is the physical basis for Schmidt et al.'s gap test: a
statistically significant gap justifies coarse-graining because the
fast modes contribute negligibly to the long-time behaviour of the
system.


% ──────────────────────────────────────────────────────────────────
\section{Spectral Coarse-Graining Theory}
\label{sec:scg_theory}
% ──────────────────────────────────────────────────────────────────

We now present the formal construction of the spectral coarse-graining
procedure following Schmidt, Caccioli, and
Aste~\cite{schmidt2024spectral}.  Let $\mathcal{G} = (V, E, w)$ be a
connected weighted graph with $N$ vertices and combinatorial Laplacian
$\mathbf{L} \in \mathbb{R}^{N \times N}$ having eigenvalues $0 =
\lambda_0 < \lambda_1 \leq \cdots \leq \lambda_{N-1}$ and
corresponding eigenvectors $\mathbf{u}_0, \mathbf{u}_1, \ldots,
\mathbf{u}_{N-1}$.

\paragraph{Step 1: Gap identification.}
For each index $k \in \{1, \ldots, N-2\}$, compute the spectral gap:
\begin{equation}
\delta_k = \lambda_{k+1} - \lambda_k.
\label{eq:spectral_gap}
\end{equation}
To determine whether $\delta_k$ is statistically anomalous, generate an
ensemble of $M = 10^4$ connected Erd\H{o}s-R\'enyi random graphs
$G_{\text{ER}}(N, p)$ with the same number of vertices and edge
probability $p = \bar{d}/(N-1)$, where $\bar{d}$ is the average degree
of $\mathcal{G}$.  For each random graph, compute the Laplacian
eigenvalues and the corresponding gap at position $k$.  The empirical
one-tailed $p$-value is:
\begin{equation}
p\text{-value}(k) = \frac{1}{M} \sum_{m=1}^{M}
\mathbf{1}\!\left[\delta_k^{(m)} \geq \delta_k^{\text{obs}}\right],
\label{eq:gap_pvalue}
\end{equation}
where $\delta_k^{(m)}$ is the gap at position $k$ in the $m$-th random
graph and $\delta_k^{\text{obs}}$ is the observed gap.  Gaps with
$p\text{-value}(k) < \alpha$ (typically $\alpha = 0.01$) are flagged
as statistically significant.

\paragraph{Step 2: Spectral clustering.}
Given a significant gap at position $k$, construct the spectral embedding
matrix:
\begin{equation}
\mathbf{U}_k = \begin{bmatrix}
\mathbf{u}_1 & \mathbf{u}_2 & \cdots & \mathbf{u}_k
\end{bmatrix} \in \mathbb{R}^{N \times k},
\label{eq:spectral_embedding}
\end{equation}
where $\mathbf{u}_1, \ldots, \mathbf{u}_k$ are the eigenvectors
corresponding to the $k$ smallest non-zero eigenvalues.  Apply $k$-means
clustering to the rows of $\mathbf{U}_k$ to partition $V$ into $k$
clusters $C_1, C_2, \ldots, C_k$~\cite{ng2002spectral,luxburg2007tutorial}.

\paragraph{Step 3: Coarse-grained Laplacian construction.}
Define the membership matrix $\mathbf{S} \in \{0,1\}^{N \times k}$
with $S_{ia} = 1$ if node $i$ belongs to cluster $C_a$.  The
coarse-grained Laplacian is:
\begin{equation}
\mathbf{L}^{(1)} = \mathbf{S}^T\,\mathbf{L}\,\mathbf{S}
\in \mathbb{R}^{k \times k}.
\label{eq:cg_laplacian}
\end{equation}
The off-diagonal entry $L^{(1)}_{ab} = -\sum_{i \in C_a, j \in C_b}
A_{ij}$ captures the total edge weight between clusters $a$ and $b$,
and the diagonal entry ensures row sums equal zero.

\paragraph{Step 4: Rescaling.}
The coarse-grained adjacency matrix $\mathbf{A}^{(1)} = \mathbf{D}^{(1)}
- \mathbf{L}^{(1)}$ is rescaled by the eigenvalue at the gap:
\begin{equation}
\mathbf{A}_R = \frac{1}{\lambda_k}\,\mathbf{A}^{(1)}.
\label{eq:rescaling}
\end{equation}
This rescaling ensures that diffusion on the coarse-grained network
approximates the slow-mode dynamics of the original network.  The
intuition is as follows: the eigenvalues of $\mathbf{L}^{(1)}$ are
approximately $\lambda_1, \ldots, \lambda_k$ scaled by the cluster
sizes; dividing by $\lambda_k$ normalises the fastest retained mode
to unit rate, compensating for the change in scale introduced by
aggregation~\cite{schmidt2024spectral}.

\paragraph{Step 5: Reconstruction error.}
The quality of the coarse-grained approximation is measured by comparing
the diffusion dynamics on the original and coarse-grained networks.
Define the reconstruction error at time $t$ as:
\begin{equation}
\varepsilon(t) = \left\|e^{-t\mathbf{L}}\,\mathbf{p}(0) -
\mathbf{S}\,e^{-t\mathbf{L}_R}\,\mathbf{S}^T\,\mathbf{p}(0)
\right\|_2^2,
\label{eq:reconstruction_error}
\end{equation}
where $\mathbf{L}_R = \mathbf{D}_R - \mathbf{A}_R$ is the Laplacian of
the rescaled coarse-grained network.  In our experiments, the
reconstruction achieves $R^2 = 1.0$ across all five topologies tested
(Section~\ref{ch:results}), indicating that the procedure preserves
diffusion dynamics with negligible error for these network
configurations.

\paragraph{Relationship to the renormalization group.}
The analogy with Wilson's renormalization
group~\cite{wilson1975renormalization} is instructive but imperfect.  In
the RG, one integrates out short-wavelength (high-frequency) modes to
obtain an effective theory for long-wavelength modes, and the coupling
constants are renormalized to preserve the partition function.  In SCG,
one integrates out fast eigenmodes (large eigenvalues) to obtain an
effective network for slow eigenmodes (small eigenvalues), and the edge
weights are rescaled to preserve diffusion dynamics.  The gap test plays
the role of identifying the ``cutoff'' scale, analogous to the momentum
cutoff in field theory.  Villegas et al.~\cite{villegas2023laplacian}
formalise this analogy more precisely through the heat kernel, though
their LRG operates on the normalised rather than combinatorial
Laplacian.

\paragraph{Why $R^2 = 1.0$?}
The perfect reconstruction observed in our experiments deserves comment.
For a network with $N = 10$ nodes that admits $k = 2$--4 clusters with a
clean spectral gap, the slow modes are exactly preserved by the
coarse-grained Laplacian (which is constructed to reproduce them), and the
fast modes contribute negligibly to the diffusion dynamics at timescales
$t \gg 1/\lambda_{k+1}$.  For larger networks where the spectral gap is
less clean or the cluster structure is fuzzier, reconstruction error
would increase.  The perfect result here reflects the small scale of the
network and the cleanness of the spectral gaps, not a general property
of the method.


% ──────────────────────────────────────────────────────────────────
\section{Graph Attention Mechanism}
\label{sec:gat_theory}
% ──────────────────────────────────────────────────────────────────

The graph attention mechanism~\cite{velickovic2018graph} computes
adaptive, data-dependent aggregation weights for message passing on
graphs.  Given a graph $\mathcal{G}$ with node feature matrix
$\mathbf{H}^{(\ell)} \in \mathbb{R}^{N \times d_\ell}$ at layer
$\ell$, the attention coefficient between nodes $i$ and $j$ is:
\begin{equation}
e_{ij} = \text{LeakyReLU}\!\left(
\mathbf{a}^T\!\left[
\mathbf{W}\,\mathbf{h}_i^{(\ell)} \;\Big\|\; \mathbf{W}\,\mathbf{h}_j^{(\ell)}
\right]\right),
\label{eq:raw_attention}
\end{equation}
where $\mathbf{W} \in \mathbb{R}^{d' \times d_\ell}$ is a shared
linear transformation, $\mathbf{a} \in \mathbb{R}^{2d'}$ is a learnable
attention vector, and $\|$ denotes concatenation.  The raw scores are
normalised via softmax over the neighbourhood:
\begin{equation}
\alpha_{ij} = \frac{\exp(e_{ij})}{\sum_{k \in \mathcal{N}(i)}
\exp(e_{ik})}.
\label{eq:softmax_attention}
\end{equation}
The updated node representation is:
\begin{equation}
\mathbf{h}_i^{(\ell+1)} = \sigma\!\left(
\sum_{j \in \mathcal{N}(i)} \alpha_{ij}\,\mathbf{W}\,
\mathbf{h}_j^{(\ell)}\right).
\label{eq:gat_update}
\end{equation}

\paragraph{Multi-head attention.}
To stabilise learning and increase expressivity, $K$ independent
attention heads are computed in parallel, and their outputs are
concatenated (in intermediate layers) or averaged (in the final
layer)~\cite{velickovic2018graph}:
\begin{equation}
\mathbf{h}_i^{(\ell+1)} = \Big\|_{k=1}^{K}\,\sigma\!\left(
\sum_{j \in \mathcal{N}(i)} \alpha_{ij}^{(k)}\,
\mathbf{W}^{(k)}\,\mathbf{h}_j^{(\ell)}\right).
\label{eq:multihead_gat}
\end{equation}
Each head has its own weight matrix $\mathbf{W}^{(k)}$ and attention
vector $\mathbf{a}^{(k)}$, learning different aspects of the
neighbourhood structure.

\paragraph{Why attention for financial networks?}
Two properties of the attention mechanism make it preferable to fixed
spectral filters (as in GCNs~\cite{kipf2017semi}) for financial
network analysis.

First, \emph{adaptivity}: the attention weights $\alpha_{ij}$ depend
on the current node features, so the aggregation pattern changes as
market conditions evolve.  In a GCN, the aggregation is determined
entirely by the adjacency matrix; if two banks have a high correlation,
their features are always aggregated with the same weight.  The
attention mechanism can learn to upweight or downweight specific
relationships depending on the input, capturing the regime-dependent
nature of financial interdependencies.

Second, \emph{interpretability}: the attention weights reveal which
inter-bank relationships the model considers most informative for
prediction.  Our experiments (Section~\ref{ch:results}) show that BNP
Paribas and HSBC receive the highest incoming attention across heads,
while Santander and Deutsche Bank exhibit the highest outgoing
attention, suggesting that peripheral banks attend primarily to core
institutions for spectral prediction.


% ──────────────────────────────────────────────────────────────────
\section{Multi-Channel Contagion Theory}
\label{sec:contagion_theory}
% ──────────────────────────────────────────────────────────────────

The agent-based model simulates contagion through three
channels~\cite{nier2007network,gai2010contagion,caccioli2015overlapping},
each governed by a distinct propagation equation.  Let
$\mathbf{E} \in \mathbb{R}^{N \times N}$ denote the bilateral exposure
matrix, $h_i(t)$ the health of institution $i$ at time step $t$ (a
composite of its CET1 ratio and liquidity coverage ratio), and
$\theta_d$ the default threshold.

\paragraph{Channel 1: Credit losses.}
When institution $j$ defaults ($h_j(t) \leq \theta_d$), creditor $i$
suffers a loss proportional to its exposure and the loss-given-default
rate:
\begin{equation}
\Delta_i^{\text{credit}}(t) = \sum_{j:\,h_j(t) \leq \theta_d}
E_{ij}\,\text{LGD},
\label{eq:credit_loss}
\end{equation}
where LGD is the loss-given-default parameter (typically $0.4$--$0.6$).
This is the classical contagion channel modelled by Nier et
al.~\cite{nier2007network} and formalised in the DebtRank
framework~\cite{battiston2012debtrank}.

\paragraph{Channel 2: Funding liquidity stress.}
Funding stress propagates through changes in perceived creditworthiness,
following the ``liquidity spiral'' mechanism of Brunnermeier and
Pedersen~\cite{brunnermeier2009market}.  The funding liquidity loss
of institution $i$ is modelled as a sigmoid function of the aggregate
distress in its neighbourhood:
\begin{equation}
\Delta_i^{\text{funding}}(t) = \gamma\;\sigma\!\left(
\beta\,\sum_{j \in \mathcal{N}(i)} (1 - h_j(t)) - \mu\right),
\label{eq:funding_loss}
\end{equation}
where $\sigma(\cdot)$ is the logistic function, $\gamma$ scales the
maximum funding loss, $\beta$ controls the sensitivity to neighbourhood
distress, and $\mu$ is a threshold parameter.  The sigmoid form captures
the non-linear nature of funding stress: small levels of neighbourhood
distress have little effect, but distress beyond a threshold triggers
rapid funding withdrawal~\cite{gai2010contagion}.

\paragraph{Channel 3: Fire-sale asset correlation.}
When distressed institutions deleverage by selling common assets, the
resulting price decline imposes mark-to-market losses on all holders of
those assets~\cite{caccioli2015overlapping}.  The fire-sale loss is:
\begin{equation}
\Delta_i^{\text{fire\text{-}sale}}(t) = \eta\,\sum_{a=1}^{A}
\omega_{ia}\,\Delta p_a(t),
\quad\text{where}\quad
\Delta p_a(t) = -\kappa\,\sum_{j:\,h_j(t) < \theta_s}
\omega_{ja}\,\text{sell}_j(t),
\label{eq:firesale_loss}
\end{equation}
where $\omega_{ia}$ is institution $i$'s holding of asset $a$,
$\Delta p_a(t)$ is the price impact of fire sales on asset $a$,
$\kappa$ is the price impact coefficient, $\theta_s$ is the stress
threshold triggering deleveraging, and $\text{sell}_j(t)$ is the volume
sold by institution $j$.  Caccioli et al.~\cite{caccioli2015overlapping}
showed that this channel can dominate direct credit losses in systems
with high leverage and concentrated portfolios.

\paragraph{Cascade termination.}
At each time step, institution $i$'s health is updated:
\begin{equation}
h_i(t+1) = h_i(t) - \Delta_i^{\text{credit}}(t) -
\Delta_i^{\text{funding}}(t) - \Delta_i^{\text{fire-sale}}(t).
\label{eq:health_update}
\end{equation}
The cascade terminates when either (a) no new defaults occur
($\forall\, i:\; h_i(t+1) > \theta_d$ if $h_i(t) > \theta_d$), or (b) a
maximum number of iterations $T_{\max}$ is reached.  In our experiments,
cascades terminate within 5--10 steps for most scenarios.

\paragraph{Ornstein-Uhlenbeck dynamics for regulatory ratios.}
Between stress scenarios, the CET1 and LCR ratios of non-defaulted
institutions are modelled as Ornstein-Uhlenbeck processes:
\begin{equation}
dX_i = \theta_{\text{OU}}\,(\mu_X - X_i)\,dt + \sigma_X\,dW_i,
\label{eq:ou_dynamics}
\end{equation}
where $X_i \in \{\text{CET1}_i, \text{LCR}_i\}$, $\theta_{\text{OU}}$
is the mean-reversion speed, $\mu_X$ is the long-run mean, and
$\sigma_X$ is the volatility.  This reflects the empirical observation
that regulatory ratios tend to revert toward target levels set by
management and supervisors, with stochastic fluctuations driven by
business conditions.


% ──────────────────────────────────────────────────────────────────
\section{Risk Indicator Formulation}
\label{sec:risk_formulation}
% ──────────────────────────────────────────────────────────────────

We propose a spectral risk indicator that combines the algebraic
connectivity and spectral radius:
\begin{equation}
\text{SCG-Risk} = 1 - \frac{\lambda_2}{\rho(\mathbf{L})},
\label{eq:scg_risk}
\end{equation}
where $\lambda_2$ is the algebraic connectivity (smallest non-zero
eigenvalue of $\mathbf{L}$) and $\rho(\mathbf{L}) =
\lambda_{N-1}$ is the spectral radius (largest eigenvalue).

\paragraph{Theoretical motivation.}
The ratio $\lambda_2/\rho$ captures the ``spectral spread'' of the
network: a network in which all non-zero eigenvalues are similar
(spread close to 1) is spectrally homogeneous, while a network with a
large gap between $\lambda_2$ and $\rho$ (spread close to 0) has a
more heterogeneous spectral structure, suggesting the presence of
bottlenecks or weakly connected communities.  Subtracting from 1 gives
a quantity that increases as the network becomes more spectrally
heterogeneous---intuitively, more ``risky'' in the sense of being
closer to fragmentation.

\paragraph{Limitations and honest assessment.}
It is essential to note that the SCG risk indicator is \emph{proposed}
rather than derived from first principles.  While $\lambda_2$ has a
rigorous interpretation as the convergence rate of diffusion processes
(Theorem~\ref{thm:convergence}), and $\rho$ bounds the maximum degree
(Theorem~\ref{thm:spectral}), the specific functional form
$1 - \lambda_2/\rho$ is a heuristic choice.

Our backtesting results (Section~\ref{ch:results}) reveal a critical
limitation: the SCG risk indicator exhibits \emph{negative} correlation
with future market volatility ($r = -0.43$ to $-0.49$), meaning that
the indicator \emph{decreases} when volatility is about to increase.
This occurs because high correlation environments (which precede high
volatility) produce dense, well-connected networks with high
$\lambda_2$ and thus low SCG risk---the opposite of what a risk
indicator should signal.  In contrast, $\lambda_2$ alone shows positive
correlation with future volatility ($r = +0.40$ to $+0.46$), and
the cross-correlation analysis reveals the underlying structure:
$\lambda_2$ vs.\ $\rho = -0.926$ and $\lambda_2$ vs.\ SCG risk $=
-0.990$.

These findings indicate that the naive SCG risk indicator captures
\emph{network fragmentation} rather than systemic stress.  A fragmented
network (low $\lambda_2$, high SCG risk) may actually be \emph{less}
systemically dangerous than a highly integrated one (high $\lambda_2$,
low SCG risk), because contagion propagates faster in
well-connected networks~\cite{acemoglu2015systemic,haldane2011systemic}.
The distinction between ``connected and vulnerable'' and ``fragmented
and safe'' is fundamental to the robust-yet-fragile property, and the
SCG risk indicator---as currently formulated---conflates connectivity
with safety rather than with vulnerability.

This is an honest negative result that informs future work: constructing
a useful spectral risk indicator requires a more nuanced treatment of
the connectivity--vulnerability relationship, possibly incorporating
the magnitude distribution of shocks (which determines whether the
system is in the ``robust'' or ``fragile'' regime) or combining spectral
properties with balance-sheet information.
